            International Mathematics Competition for University Students
                          July 25–30 2009, Budapest, Hungary
                                                               Day 1

Problem 1.
Suppose that f and g are real-valued functions on the real line and f (r) ≤ g(r) for every rational r. Does this
imply that f (x) ≤ g(x) for every real x if
    a) f and g are non-decreasing?
    b) f and g are continuous?
                                                                                            √            √
Solution. a) No. Counter-example: f and g can be chosen as the characteristic functions of [ 3, ∞) and ( 3, ∞),
respectively.
    b) Yes. By the assumptions g − f is continuous on the whole real line and nonnegative on the rationals. Since
any real number can be obtained as a limit of rational numbers we get that g − f is nonnegative on the whole real
line.

Problem 2.
Let A, B and C be real square matrices of the same size, and suppose that A is invertible. Prove that if (A−B)C =
BA−1 , then C(A − B) = A−1 B.
Solution. A straightforward calculation shows that (A−B)C = BA−1 is equivalent to AC −BC −BA−1 +AA−1 =
I, where I denotes the identity matrix. This is equivalent to (A − B)(C + A−1 ) = I. Hence, (A − B)−1 = C + A−1 ,
meaning that (C + A−1 )(A − B) = I also holds. Expansion yields the desired result.

Problem 3.
In a town every two residents who are not friends have a friend in common, and no one is a friend of everyone else.
Let
Pn us 2number   the residents from 1 to n and let ai be the number of friends of the i-th resident. Suppose that
      a  = n 2 − n. Let k be the smallest number of residents (at least three) who can be seated at a round table
   i=1 i
in such a way that any two neighbors are friends. Determine all possible values of k.
Solution. Let us define the simple, undirected graph G so that the vertices of G are the town’s residents and the
edges of G are the friendships between the residents. Let V (G) = {v1 , v2 , . . . , vn } denote the vertices of G; ai is
degree of vi for every i. Let E(G) denote the edges of G. In this terminology, the problem asks us to describe the
length k of the shortest cycle in G.
    Let us count the walks of length 2 in G, that is, the ordered triples (vi , vj , vl ) of vertices with         P vi vj , vj vl ∈ E(G)
(i = l being allowed). For a given j the number is obviously a2j , therefore the total number is ni=1 a2i = n2 − n.
    Now we show that there is an injection f from the set of ordered pairs of distinct vertices to the set of these
walks. For vi vj ∈    / E(G), let f (vi , vj ) = (vi , vl , vj ) with arbitrary l such that vi vl , vl vj ∈ E(G). For vi vj ∈ E(G), let
f (vi , vj ) = (vi , vj , vi ). f is an injection since for i 6= l, (vi , vj , vl ) can only be the image of (vi , vl ), and for i = l, it
can only be the image of (vi , vj ).                                                     P
    Since the number of ordered pairs of distinct vertices is n2 − n, ni=1 a2i ≥ n2 − n. Equality holds iff f is
surjective, that is, iff there is exactly one l with vi vl , vl vj ∈ E(G) for every i, j with vi vj ∈          / E(G) and there is no
such l for any i, j with vi vj ∈ E(G). In other words, iff G contains neither C3 nor C4 (cycles of length 3 or 4), that
is, G is either a forest (a cycle-free graph) or the length of its shortest cycle is at least 5.
    It is easy to check that if every two vertices of a forest are connected by a path of length at most 2, then the
forest is a star (one vertex is connected to all others by an edge). But G has n vertices, and none of them has
degree n − 1. Hence G is not forest, so it has cycles. On the other hand, if the length of a cycle C of G is at
least 6 then it has two vertices such that both arcs of C connecting them are longer than 2. Hence there is a path
connecting them that is shorter than both arcs. Replacing one of the arcs by this path, we have a closed walk
shorter than C. Therefore length of the shortest cycle is 5.
    Finally, we must note that there is at least one G with the prescribed properties – e.g. the cycle C5 itself
satisfies the conditions. Thus 5 is the sole possible value of k.




                                                                    1
Problem 4.
Let p(z) = a0 + a1 z + a2 z 2 + · · · + an z n be a complex polynomial. Suppose that 1 = c0 ≥ c1 ≥ · · · ≥ cn ≥ 0 is a
sequence of real numbers which is convex (i.e. 2ck ≤ ck−1 + ck+1 for every k = 1, 2, . . . , n − 1), and consider the
polynomial
                                       q(z) = c0 a0 + c1 a1 z + c2 a2 z 2 + · · · + cn an z n .
   Prove that
                                                          max q(z) ≤ max p(z) .
                                                          |z|≤1                   |z|≤1

Solution. The polynomials p and q are regular on the complex plane, so by the Maximum Principle, max|z|≤1 |q(z)| =
max|z|=1 |q(z)|, and similarly for p. Let us denote Mf = max|z|=1 |f (z)| for any regular function f . Thus it suffices
to prove that Mq ≤ Mp .
    First, note that we can assume cn = 0. Indeed,
                                               P       for cn = 1, we get p = q and the statement is trivial; otherwise,
                                                     cj −cn                             cj −cn
q(z) = cn p(z) + (1 − cn )r(z), where r(z) = nj=0 1−c     n
                                                            aj z j . The sequence c′j = 1−c  n
                                                                                               also satisfies the prescribed
conditions (it is a positive linear transform of the sequence cn with c0 = 1), but c′n = 0 too, so we get Mr ≤ Mp .
                                                                             ′

This is enough: Mq = |q(z0 )| ≤ cn |p(z0 )| + (1 − cn )|r(z0 )| ≤ cn Mp + (1 − cn )Mr ≤ Mp .
    Using the Cauchy formulas, we can express the coefficients aj of p from its values taken over the positively
oriented circle S = {|z| = 1}:                    Z                     Z
                                                1      p(z)           1    p(z)
                                         aj =               dz   =              |dz|
                                               2πi S z j+1           2π S z j
for 0 ≤ j ≤ n, otherwise                                      Z
                                                                   p(z)
                                                                      j
                                                                        |dz| = 0.
                                                                  S z
   Let us use these identities to get a new formula for q, using only the values of p over S:
                                                  n
                                                  X     Z              
                                                                  −j
                                      2π · q(w) =    cj      p(z)z |dz| wj .
                                                                  j=0             S


We can exchange the order of the summation and the integration (sufficient conditions to do this obviously apply):
                                                                
                                                Z   n
                                                    X
                                    2π · q(w) =       cj (w/z)j  p(z)|dz|.
                                                                  S      j=0

It would be nice if the integration kernel (the sum between the brackets) was real. But this is easily arranged – for
−n ≤ j ≤ −1, we can add the conjugate expressions, because by the above remarks, they are zero anyway:
                                   Xn     Z                   Xn        Z            
                                                   −j       j                      −j
                      2π · q(w) =      cj     p(z)z |dz| w =         c|j|    p(z)z |dz| wj ,
                                    j=0               S                                   j=−n           S

                                                                             
                                          Z           n
                                                      X                                            Z
                                                                           j
                           2π · q(w) =                     c|j| (w/z)            p(z)|dz| =            K(w/z)p(z)|dz|,
                                              S   j=−n                                              S

where
                                                           n
                                                           X                                n
                                                                                            X
                                                                          j
                                          K(u) =                      c|j| u = c0 + 2             cj R(uj )
                                                          j=−n                              j=1

for u ∈ S.                                                                                 R
  R Let us examine K(u).   It is a real-valued function. Again from  the  Cauchy formulas,   S K(u)|du| = 2πc0 = 2π.
If S |K(u)||du| = 2π still holds (taking the absolute value does not increase the integral), then for every w:
                          Z                        Z                             Z
              2π|q(w)| =      K(w/z)p(z)|dz| ≤        |K(w/z)| · |p(z)||dz| ≤ Mp    |K(u)||du| = 2πMp ;
                             S                                    S                                           S
                                                                                      R                       R
this would conclude the proof. So it suffices to prove that                               S |K(u)||du| =          S K(u)|du|, which is to say, K is
non-negative.

                                                                              2
                               c_{|j|}


                                d_1 F_1



                                     d_2 F_2



                                                           j

   Now let us decompose K into a sum using the given conditions for the numbers cj (including cn = 0). Let
                                                                                                   Pk−1
dk = ck−1 − 2ckP+ ck+1 for k = 1, . . . , n (setting cn+1 = 0); we know that dk ≥ 0. Let Fk (u) = j=−k+1    (k − |j|)uj .
                   n
Then K(u) = k=1 dk Fk (u) by easy induction (or see Figure for a graphical illustration). So it suffices to prove
that Fk (u) is real and Fk (u) ≥ 0 for u ∈ S. This is reasonably well-known (as Fkk is the Fejér kernel), and also very
easy:
                         Fk (u) = (1 + u + u2 + · · · + uk−1 )(1 + u−1 + u−2 + · · · + u−(k−1) ) =
               = (1 + u + u2 + · · · + uk−1 )(1 + u + u2 + · · · + uk−1 ) = |1 + u + u2 + · · · + uk−1 |2 ≥ 0
   This completes the proof.

Problem 5.
Let n be a positive integer. An n-simplex in Rn is given by n + 1 points P0 , P1 , . . . , Pn , called its vertices, which
do not all belong to the same hyperplane. For every n-simplex S we denote by v(S) the volume of S, and we write
C(S) for the center of the unique sphere containing all the vertices of S.
    Suppose that P is a point inside an n-simplex S. Let Si be the n-simplex obtained from S by replacing its i-th
vertex by P . Prove that

                             v(S0 )C(S0 ) + v(S1 )C(S1 ) + · · · + v(Sn )C(Sn ) = v(S)C(S).

Solution 1. We will prove this by induction on n, starting with n = 1. In that case we are given an interval [a, b]
with a point p ∈ (a, b), and we have to verify

                                                    b+p           p+a           b+a
                                          (b − p)       + (p − a)     = (b − a)     ,
                                                     2             2             2
which is true.
   Now let assume the result is true for n − 1 and prove it for n. We have to show that the point
                                                               n
                                                               X v(Sj )
                                                        X=                  O(Sj )
                                                                     v(S)
                                                               j=0

has the same distance to all the points P0 , P1 , . . . , Pn .              Let i ∈ {0, 1, 2, . . . , n} and define the sets
Mi = {P0 , P1 , . . . , Pi−1 , Pi+1 , . . . , Pn }. The set of all points having the same distance to all points in Mi is a
line hi orthogonal to the hyperplane Ei determined by the points in Mi . We are going to show that X lies on every
hi . To do so, fix some index i and notice that

                                         v(Si )          v(S) − v(Si ) X        v(Sj )
                               X=               O(Si ) +              ·                    O(Sj )
                                         v(S)               v(S)             v(S) − v(Si )
                                                                        j6=i
                                                                        |          {z          }
                                                                                     Y

and O(Si ) lies on hi , so that it is enough to show that Y lies on hi .
    A map f : R>0 → Rn will be called affine if there are points A, B ∈ Rn such that f (λ) = λA + (1 − λ)B.
Consider the ray g starting in Pi and passing through P . For λ > 0 let Pλ = (1 − λ)P + λPi , so that Pλ is an
affine function describing the points of g. For every such λ let Sjλ be the n-simplex obtained from S by replacing
the j-th vertex by Pλ . The point O(Sjλ ) is the intersection of the fixed line hj with the hyperplane orthogonal to

                                                                     3
g and passing through the midpoint of the segment Pi Pλ which is given by an affine function. This implies that
                                                     v(Sj )
also O(Sjλ ) is an affine function. We write ϕj = v(S)−s(S  i)
                                                               , and then
                                                        X
                                                 Yλ =         ϕj O(Sjλ )
                                                              j6=i

is an affine function. We want to show that Yλ ∈ hi for all λ (then specializing to λ = 1 gives the desired result).
It is enough to do this for two different values of λ.
     Let g intersect the sphere containing the vertices of S in a point Z; then Z = Pλ for a suitable λ > 0, and we
have O(Sjλ ) = O(S) for all j, so that Yλ = O(S) ∈ hi . Now let g intersect the hyperplane Ei in a point Q; then
Q = Pλ for some λ > 0, and Q is different from Z. Define T to be the (n − 1)-simplex with vertex set Mi , and
let Tj be the (n − 1)-simplex obtained from T by replacing the vertex Pj by Q. If we write v ′ for the volume of
(n − 1)-simplices in the hyperplane Ei , then

                                      v ′ (Tj )   v(Sjλ )     v(Sjλ )
                                                =         = P
                                      v ′ (T )    v(S)                λ
                                                             k6=i v(Sk )
                                                      λv(Sj )           v(Sj )
                                              =P                   =               = ϕj .
                                                     k6=i λv(S k )   v(S) − v(Si )
                                                                                                     P
If p denotes the orthogonal projection onto Ei then p(O(Sjλ )) = O(Tj ), so that p(Yλ ) =              j6=i ϕj O(Tj ) equals O(T )
by induction hypothesis, which implies Yλ ∈ p−1 (O(T )) = hi , and we are done.
Solution 2. For n = 1, the statement is checked easily.
   Assume n ≥ 2. Denote O(Sj ) − O(S) by qj and Pj − P by pj . For all distinct j and k in the range 0, ..., n the
point O(Sj ) lies on a hyperplane orthogonal to pk and Pj lies on a hyperplane orthogonal to qk . So we have
                                               (
                                                 hpi , qj − qk i = 0
                                                 hqi , pj − pk i = 0

for all j 6= i 6= k. This means that the value hpi , qj i is independent of j as long as j 6= i, denote this value by λi .
Similarly, hqi , pj i = µi for some µi . Since n ≥ 2, these equalities imply that all the λi and µi values are equal, in
particular, hpi , qj i = hpj , qi i for any i and j.
    We claim that for such pi and qi , the volumes
                                            Vj = | det(p0 , ..., pj−1 , pj+1 , ..., pn )|
and
                                            Wj = | det(q0 , ..., qj−1 , qj+1 , ..., qn )|
are proportional. Indeed, first assume that p0 , ..., pn−1 and q0 , ..., qn−1 are bases of Rn , then we have
                                          1                            
                         Vj =                          det hpk , ql i k6=j =
                               | det(q0 , ..., qn−1 )|                     l<n
                                          1                                     det(p0 , ..., pn−1 )
                             =                         det((hpk , ql i)) l6=j =                      Wj .
                               | det(q0 , ..., qn−1 )|                  k<n     det(q0 , ..., qn−1 )
If our assumption did not hold after any reindexing of the vectors pi and qi , then both pi and qi span a subspace
of dimension at most n − 1 P  and all the volumes are 0.
    Finally, it is clear that   qj Wj / det(q0 , ..., qn ) = 0: the weight of pj is the height of 0 over the hyperplane
spanned by the rest of the vectors qk relative to the height of pj over the same hyperplane, so the sum is parallel
to all the faces of the simplex spanned by q0 , ..., qn . By the argument above, we can change the weights to the
proportional set of weights Vj / det(p0 , ..., pn ) and the sum will still be 0. That is,
                    X        Vj              X                v(Sj )
               0=       qj                 =  (O(Sj ) − O(S))        =
                        det(p0 , ..., pn )                    v(S)
                     1 X                        X               1 X                        
                 =          O(Sj )v(Sj ) − O(S)       v(Sj ) =         O(Sj )v(Sj ) − O(S)v(S) ,
                   v(S)                                         v(S)
q.e.d.

                                                                 4
